\section{Comparison with other input-level detectors}
\label{sec:robustness:detectors}\label{sec:detectors}

We ported \DetectorsComparedCount{} input-level detectors to exactly the splits, the threshold rule and the false-positive budgets PSBD is scored on, each checked numerically against its reference implementation, and ran all of them on all \DetectorsComparedModels{} backdoored models. \Cref{tab:detectors-summary} gives the mean per dataset and \cref{tab:detectors-auroc} in \cref{app:detectors} every attack. No single defense wins everywhere. PSBD-TM leads on \DetectorsAurocDatasetsOursLeads{}, and \DetectorsAurocDatasetsOthersLead{}.

% generated by scripts/paper/tab_detectors.py from results/coverage/coverage.json (54 of 56 successful cells carry every defense), results/*/detectors/*_metrics.json, results/*/psbd_metrics.json at 2026-09-29T18:05:43.231535+00:00, commit b2d32cf11708d5de965d3e13604c863a3ad9b493-dirty. Do not edit by hand.
\begin{table}[htbp]
\centering
\small
\caption{Mean AUROC of every defense per dataset over the 54 backdoored ViT-B/16 models that carry all of them, ours first. Best in each column is bold and second best underlined, and gray marks a mean below chance. PSBD-RD is our adaptation of the original PSBD site. \Cref{tab:detectors-auroc} gives every attack and poison rate.}
\label{tab:detectors-summary}
\begin{adjustbox}{max width=\linewidth}
\begin{tabular}{lrrrrr}
\toprule
defense & CIFAR-10 & CIFAR-100 & GTSRB & Tiny & all \\
\midrule
PSBD-TM & \underline{0.919} & \textbf{0.979} & 0.984 & \textbf{0.977} & \textbf{0.963} \\
PSBD-RD & 0.831 & 0.890 & 0.855 & \underline{0.965} & 0.885 \\
Conf & 0.593 & 0.813 & \textcolor{black!45}{0.480} & 0.877 & 0.688 \\
STRIP & 0.850 & 0.876 & 0.878 & 0.831 & 0.857 \\
Scale-Up & 0.784 & 0.758 & 0.639 & 0.723 & 0.728 \\
Scale-Up$^\dagger$ & 0.563 & 0.641 & 0.817 & 0.544 & 0.636 \\
IBD-PSC & 0.639 & 0.833 & 0.531 & 0.931 & 0.732 \\
IBD-PSC$^\ast$ & 0.917 & 0.943 & 0.942 & 0.944 & \underline{0.936} \\
TeCo & 0.677 & 0.808 & 0.655 & 0.841 & 0.743 \\
CD-L & 0.869 & 0.810 & 0.839 & 0.711 & 0.808 \\
Beatrix & \textbf{0.977} & 0.965 & \textbf{0.998} & 0.655 & 0.896 \\
TED & 0.861 & \underline{0.976} & \underline{0.993} & 0.733 & 0.885 \\
SentiNet & \textcolor{black!45}{0.326} & 0.762 & \textcolor{black!45}{0.217} & \textcolor{black!45}{0.408} & \textcolor{black!45}{0.418} \\
\bottomrule
\end{tabular}
\end{adjustbox}
\end{table}


The placement is what makes the method competitive. PSBD-RD, our adaptation of the original site, ranks 5th of \DetectorsAurocDefensesRanked{} defenses by mean AUROC, while PSBD-TM ranks 1st by AUROC and by the true-positive rate at both false-positive budgets. The strongest competitor is \DetectorsAurocBestCompetitorName{}, which reads \DetectorsAurocBestCompetitor{} mean AUROC against \DetectorsAurocOurs{} for PSBD-TM. The paired margin is \DetectorsAurocMargin{} [\DetectorsAurocMarginLow{}, \DetectorsAurocMarginHigh{}] in AUROC, \DetectorsTprOneZeroMargin{} [\DetectorsTprOneZeroMarginLow{}, \DetectorsTprOneZeroMarginHigh{}] in TPR at a 10\% budget and \DetectorsTprTwoZeroMargin{} [\DetectorsTprTwoZeroMarginLow{}, \DetectorsTprTwoZeroMarginHigh{}] at 20\%. Every interval contains zero, so PSBD-TM leads that port on all three measures without a lead we can separate from noise, and it leads every other defense clearly.

We count only the calibrated IBD-PSC port as the competitor, because the port at the paper's fixed amplification factor saturates on a ViT and reads \DetectorsAurocIbdPsc{}. SentiNet reads below chance, which \cref{app:detectors} discusses.

\section{An adaptive attacker and the union of probes}
\label{sec:robustness:adaptive}

The PSBD paper reports that its rate rule survives an attacker who lowers the prediction shift of poisoned samples. We trained a stronger attacker on ViT. It adds a hinge loss that pushes the poisoned samples' PSU up to the clean samples' under PSBD-TM, the exact probe the defender uses. The attacker optimizes this loss jointly with the backdoor on CIFAR-100 for BadNets, Blend, BPP, LF and Adaptive-Blend at three poison rates. \AdaptiveVitCells{} evasive models keep an ASR of at least \AdaptiveHighAsr{} and lose \AdaptiveVitCaCost{} clean accuracy on average (\cref{tab:adaptive} in \cref{app:adaptive}).

Against a single fixed probe the attacker wins. PSBD-TM falls from a mean AUROC of \AdaptiveVitProbedBase{} to \AdaptiveVitProbedEvade{}, with the rate rule already inside these numbers. The evasion does not transfer between probes, though, because the attacker's gradient only touches the features its one probe reads, and the two probes it never saw still read \AdaptiveVitTransfer{}.

The defense we recommend is therefore a union of probes. Let $r_j(x)$ be the rank of input $x$'s PSU under probe $j$ among the clean validation inputs. The union score is $u(x) = \min_j r_j(x)$ and $x$ is flagged when $u(x)$ falls below the \HeadlineQuantilePercent{} percentile of $u$ over the clean validation set, which keeps the false-positive rate at the budget whatever the probes' correlation. A union of PSBD-TM, dropout at the attention input and gain scaling of the MLP LayerNorm output restores a mean AUROC of \AdaptiveVitUnion{} on the evasive models and flags \AdaptiveVitUnionTpr{} of triggered inputs at a measured false-positive rate of \AdaptiveVitUnionFpr{}. Dropping the attacked probe, as if the defender knew which it was, reads only \AdaptiveVitOracleGap{} higher, so the defender does not need to know. Evading one probe already costs the attacker up to \AdaptiveVitCaCostMaxPoints{} clean-accuracy points, and an attacker trained against all three is untested. On the ordinary models, pairing PSBD-TM with token masking on the attention branch output gains \ProbeUnionTmBranchGain{}, because placements disagree on which inputs are fragile (\cref{app:union}).
